\subsection{ニュートン法型}

少々離れたところから話を始める.
ある関数$f(x)$について,\,$f(a)$が既知の時,その近傍の点$x=a+h$で
\[
  f(a+h) \approx f(a)+f'(a)(x-a)
\]
が成り立つ.
つまり,\,$f(a)$の値を元に,\,$f(a+h)$を近似したのである.
方程式
\[
  f(x)=x^2-2=0
\]
の正の解 $\sqrt{2}$ を求めたいとする.いまの近似値を $x=a_n$ とし,
点 $(a_n,f(a_n))$ における接線を引く.その接線の傾きは
$f'(a_n)=2a_n$ だから,接線の方程式は
\[
  y=f(a_n)+2a_n(x-a_n)
\]
である.\,$y=0$ と置いて $x$ 軸との交点を求めると,
\begin{align*}
  a_{n+1}
    &=a_n-\frac{f(a_n)}{f'(a_n)}\\
    &=a_n-\frac{a_n^2-2}{2a_n}\\
    &=\frac{a_n^2+2}{2a_n}
\end{align*}
これが
\[
  a_{n+1}=\frac{a_n^2+2}{2a_n}
\]
の意味である.

例えば $a_1=2$ とすると,
\[
  a_2=\frac{3}{2},
  \quad
  a_3=\frac{17}{12},
  \quad
  a_4=\frac{577}{408}
\]
となり,これらは $\sqrt{2}$ に急速に近づく.この速さは,誤差を
$e_n=a_n-\sqrt{2}$ と置くと見える.
\begin{align*}
  e_{n+1}
    &=\frac{a_n^2+2}{2a_n}-\sqrt{2}\\
    &=\frac{(a_n-\sqrt{2})^2}{2a_n}\\
    &=\frac{e_n^2}{2a_n}
\end{align*}
誤差がほぼ二乗されるため,正しい近似に近づくと桁数が一気に増える.

この漸化式の解き方は,一般項を閉じた式にすることではない.
「接線を使って,次の近似値をどう設計するか」を考え,その設計が
どのような速さで誤差を減らすかを調べている.付録では,この考えを
漸化式の収束,固定点,誤差評価,微分方程式の数値解法へ広げる.
